% Part: turing-machines 
% Chapter: undecidability 
% Section: verification

\documentclass[../../../include/open-logic-section]{subfiles}

\begin{document}

\olfileid[pt]{tur}{und}{ver} 
\olsection{Verificando a Representação}

\begin{explain}
A fim de verificar que nossa representação funciona, temos de provar
duas coisas. Primeiro, temos de mostrar que se $M$ para na entrada~$w$,
então $!T(M, w) \lif !E(M, w)$ é válida. Em seguida, temos de mostrar a
recíproca, isto é, que se $!T(M, w) \lif !E(M, w)$ é válida, então $M$
de fato eventualmente para quando executada na entrada~$w$.

A estratégia para prová-las é bem diferente. Para o primeiro
resultado, temos de mostrar que !!a{sentence} da lógica de primeira
ordem (a saber, $!T(M, w) \lif !E(M, w)$) é válida. A maneira mais
fácil de fazer isso é dar !!a{derivation}. Nossa prova, porém, deve
funcionar para todos os $M$ e $w$, de modo que não há realmente uma
única !!{sentence} para a qual tenhamos de dar !!a{derivation}, mas
infinitas. Assim, o melhor que podemos fazer é provar por indução que,
sejam quais forem $M$ e~$w$, e por mais passos que $M$ leve para parar
na entrada~$w$, haverá !!a{derivation} de $!T(M, w) \lif !E(M, w)$.

Naturalmente, nossa indução procederá sobre o número de passos que $M$
leva antes de atingir uma configuração de parada. Em nossa prova
indutiva, estabeleceremos que para cada passo~$n$ da execução de $M$ na
entrada~$w$, $!T(M, w) \Entails !C(M, w, n)$, em que $!C(M, w, n)$
descreve corretamente a configuração de $M$ executada em~$w$ após~$n$
passos. Ora, se $M$ para na entrada~$w$ após, digamos, $n$ passos,
$!C(M, w, n)$ descreverá uma configuração de parada. Mostraremos também
que $!C(M, w, n) \Entails !E(M, w)$ sempre que $!C(M, w, n)$ descreve
uma configuração de parada. Assim, se $M$ para na entrada~$w$, então
para algum~$n$, $M$ estará em uma configuração de parada após $n$~
passos. Logo, $!T(M, w) \Entails !C(M, w, n)$ em que $!C(M, w, n)$
descreve uma configuração de parada, e como nesse caso $!C(M, w, n)
\Entails !E(M, w)$, obtemos que $T(M, w) \Entails !E(M, w)$, isto é,
que $\Entails !T(M, w) \lif !E(M, w)$.

A estratégia para a recíproca é bem diferente. Aqui supomos que
$\Entails !T(M, w) \lif !E(M, w)$ e temos de provar que $M$ para na
entrada~$w$. Da hipótese obtemos que $!T(M, w) \Entails !E(M, w)$, isto
é, $!E(M, w)$ é verdadeira em toda !!{structure} em que $!T(M, w)$ é
verdadeira. Assim, descreveremos !!a{structure}~$\Struct{M}$ em que
$!T(M, w)$ é verdadeira: seu domínio será $\Nat$, e a interpretação de
todos os $\Obj Q_q$ e $\Obj S_\sigma$ será dada pelas configurações
de~$M$ durante uma execução na entrada~$w$. Assim, por exemplo,
$\Sat{M}{\Obj Q_q(\num{m}, \num{n})}$ se, e somente se, $T$, quando
executada na entrada~$w$ por $n$ passos, está no estado~$q$ e
examinando o quadrado~$m$. Ora, como $!T(M, w) \Entails !E(M, w)$ por
hipótese, e como $\Sat{M}{!T(M, w)}$ por construção, $\Sat{M}{!E(M,
w)}$. Mas $\Sat{M}{!E(M, w)}$ se, e somente se, há algum $n \in
\Domain{M} = \Nat$ tal que $M$, executada na entrada~$w$, está em uma
configuração de parada após~$n$ passos.
\end{explain}

\begin{defn} 
Seja $!C(M, w, n)$ a !!{sentence}
\[ 
\Obj Q_q(\num{m}, \num{n}) \land \Obj S_{\sigma_0}(\num{0}, \num{n})
\land \dots \land \Obj S_{\sigma_k}(\num{k}, \num{n}) \land
\lforall[x][(\num{k} < x \lif \Obj S_\TMblank(x, \num{n}))]
\] 
em que $q$ é o estado de $M$ no instante~$n$, $M$ está examinando o
quadrado~$m$ no instante~$n$, o quadrado~$i$ contém o símbolo~$\sigma_i$
no instante~$n$ para $0 \le i \le k$, e $k$ é o quadrado não-vazio
mais à direita da fita no instante~$0$, ou o quadrado mais à direita
que o cabeçote de fita visitou após $n$ passos, o que for maior.
\end{defn}

\begin{lem}\ollabel{lem:halt-config-implies-halt}
Se $M$ executada na entrada~$w$ está em uma configuração de parada após
$n$ passos, então $!C(M, w, n) \Entails !E(M, w)$.
\end{lem}

\begin{proof}
Suponha que $M$ pare para a entrada~$w$ após $n$ passos.
Há algum estado~$q$, quadrado~$m$ e símbolo~$\sigma$ tais que:
\begin{enumerate} 
\item Após $n$ passos, $M$ está no estado~$q$ examinando o
  quadrado~$m$ no qual aparece~$\sigma$.
\item A função de transição $\delta(q, \sigma)$ é indefinida.
\end{enumerate}
$!C(M, w, n)$ é a descrição dessa configuração e incluirá as cláusulas
$\Obj Q_{q}(\num{m}, \num{n})$ e $\Obj S_{\sigma}(\num{m}, \num{n})$.
Essas cláusulas juntas implicam $!E(M, w)$:
\[
\lexists[x][\lexists[y][(\bigvee_{\tuple{q, \sigma} \in
      X}(\Obj Q_q(x, y) \land \Obj S_\sigma(x, y)))]]
\]
já que $\Obj Q_{q'}(\num{m}, \num{n}) \land S_{\sigma'}(\num{m},
\num{n}) \Entails \bigvee_{\tuple{q, \sigma} \in X} (\Obj Q_q(\num{m},
\num{n}) \land \Obj S_{\sigma}(\num{m}, \num{n}))$, pois
$\tuple{q',\sigma'} \in X$. 
\end{proof}

\begin{explain}
Assim, se $M$ para na entrada $w$, então há algum~$n$ tal que $!C(M,
w, n) \Entails !E(M,w)$.  Mostraremos agora que para qualquer
instante~$n$, $!T(M, w) \Entails !C(M, w, n)$.
\end{explain}

\begin{lem}
\ollabel{lem:config}
Para cada $n$, se $M$ não parou após $n$ passos, $!T(M, w)
\Entails !C(M, w, n)$.
\end{lem}

\begin{proof}
Base da indução: Se $n = 0$, então as fórmulas componentes de $!C(M,
w, 0)$ são também fórmulas componentes de $!T(M, w)$, logo implicadas
por ela.

Hipótese de indução: Se $M$ não parou antes do $n$-ésimo passo,
então $!T(M,w) \Entails !C(M, w, n)$. Temos de mostrar que (a menos
que $!C(M, w, n)$ descreva uma configuração de parada), $!T(M, w)
\Entails !C(M, w, n+1)$.

Suponha $n > 0$ e que, após $n$ passos, $M$ iniciada em $w$ esteja no
estado~$q$ examinando o quadrado~$m$. Como $M$ não para após~$n$
passos, deve haver no programa de~$M$ uma instrução de uma das três
formas seguintes:

\begin{enumerate} 
\item \ollabel{right} $\delta(q, \sigma) = \tuple{q', \sigma', \TMright}$

\item \ollabel{left} $\delta(q, \sigma) = \tuple{q', \sigma', \TMleft}$

\item \ollabel{stay} $\delta(q, \sigma) = \tuple{q', \sigma', \TMstay}$
\end{enumerate}

Consideraremos cada um desses três casos por sua vez. 

\begin{enumerate} 
\item Suponha que haja uma instrução da forma~\olref{right}.
  Por \olref[rep]{defn:tm-descr}\olref[rep]{rep-right}, isso significa que
\begin{align*} 
& \lforall[x][\lforall[y][((\Obj Q_{q}(x,
  y) \land \Obj S_\sigma(x, y)) \lif {}]]\\
& \qquad (\Obj Q_{q'}(x',
  y') \land \Obj S_{\sigma'}(x, y') \land !A(x, y)))  
\intertext{é uma fórmula componente de $!T(M,w)$. Isso implica a
  seguinte !!{sentence} (instanciação universal, $\num{m}$ por~$x$ e
  $\num{n}$ por~$y$):}
& (\Obj Q_{q}(\num{m}, \num{n}) \land \Obj S_{\sigma}(\num{m},
\num{n})) \lif {}\\
& \qquad (\Obj Q_{q'}(\num{m}', \num{n}') \land
\Obj S_{\sigma'}(\num{m}, \num{n}') \land !A(\num{m}, \num{n})).
\intertext{Por hipótese de indução, $!T(M, w) \Entails !C(M, w, n)$,
  isto é,}
& \Obj Q_q(\num{m}, \num{n}) \land \Obj S_{\sigma_0}(\num{0}, \num{n})
\land \dots \land \Obj S_{\sigma_k}(\num{k}, \num{n}) \land \\
& \qquad\lforall[x][(\num{k} < x \lif \Obj S_\TMblank(x, \num{n}))]\\
\intertext{Como após $n$ passos o quadrado de fita~$m$ contém~$\sigma$,
  a fórmula componente correspondente é~$\Obj S_\sigma(\num{m},
  \num{n})$, de modo que isso implica:}
& \Obj Q_{q}(\num{m}, \num{n}) \land \Obj S_{\sigma}(\num{m},
   \num{n})
\intertext{Obtemos agora}
& \Obj Q_{q'}(\num{m}', \num{n}') \land \Obj S_{\sigma'}(\num{m},
  \num{n}') \land {}\\
& \qquad\Obj S_{\sigma_0}(\num{0}, \num{n}') \land \dots \land
  \Obj S_{\sigma_k}(\num{k}, \num{n}') \land {}\\
& \qquad \lforall[x][(\num{k} < x \lif \Obj S_\TMblank(x, \num{n}'))]
\end{align*}
como segue: A primeira linha vem diretamente do consequente do
condicional precedente, por modus ponens. Cada fórmula componente na
linha do meio---que exclui $S_{\sigma_m}(\num{m},\num{n}')$---decorre
da fórmula componente correspondente em~$!C(M, w, n)$ junto com
$!A(\num{m}, \num{n})$.

Se $m < k$, $!T(M,w) \Proves \num{m} < \num {k}$
(\olref[rep]{prop:mlessk}) e, por transitividade de~$<$, temos
$\lforall[x][(\num{k} < x \lif \num{m} < x)]$.  Se $m = k$, então
$\lforall[x][(\num{k} < x \lif \num{m} < x)]$ apenas por lógica.  A
última linha decorre então da fórmula componente correspondente em
$!C(M, w, n)$, de $\lforall[x][(\num{k} < x \lif \num{m} < x)]$ e de
$!A(\num{m}, \num{n})$.  Se $m<k$, isto já é $!C(M, w, n+1)$.

Agora suponha $m=k$. Nesse caso, após $n+1$ passos, o cabeçote de
fita também visitou o quadrado~$k+1$, que agora é o quadrado mais à
direita visitado.  Assim, $!C(M, w, n+1)$ tem uma nova fórmula
componente, $\Obj S_\TMblank(\num{k}',\num{n}')$, e a última fórmula
componente é $\lforall[x][(\num{k}' < x \lif \Obj S_\TMblank(x,
\num{n}'))]$. Temos de verificar que essas duas !!{sentence}s também
são implicadas.

Já temos $\lforall[x][(\num{k} < x \lif \Obj S_\TMblank(x,
  \num{n}'))]$. Em particular, isso nos dá $\num{k} < \num{k}' \lif
\Obj S_\TMblank(\num{k}', \num{n}')$. Do axioma $\lforall[x][x <
  x']$ obtemos $\num{k} < \num{k}'$. Por modus ponens, segue-se $\Obj
S_\TMblank(\num{k}',\num{n}')$.

Além disso, como $!T(M,w) \Proves \num{k} < \num{k}'$, o axioma da
transitividade de~$<$ nos dá $\lforall[x][(\num{k}' < x \lif \Obj
  S_\TMblank(x, \num{n}'))]$. (Deixamos a verificação disso como
exercício.)

\item Suponha que haja uma instrução da forma~\olref{left}.
  Então, por \olref[rep]{defn:tm-descr}\olref[rep]{rep-left},
\begin{align*} 
& \lforall[x][\lforall[y][((\Obj Q_{q}(x', y) \land \Obj
    S_{\sigma}(x', y)) \lif {}]]\\
& \qquad (\Obj Q_{q'}(x, y') \land \Obj
  S_{\sigma'}(x', y') \land !A(x, y))) \land {}\\
& \lforall[y][((\Obj Q_{q_i}(\num{0}, y) \land \Obj S_{\sigma}(\num{0},
    y)) \lif {}]\\
& \qquad (\Obj Q_{q_j}(\num{0}, y') \land \Obj S_{\sigma'}(\num{0}, y')
  \land !A(\num{0}, y)))
\intertext{é uma fórmula componente de $!T(M,w)$. Se $m>0$, então seja
  $l = m - 1$ (isto é, $m = l+1$). A primeira fórmula componente da
  !!{sentence} acima implica o seguinte:}
& (\Obj Q_{q}(\num{l}', \num{n}) \land \Obj S_{\sigma}(\num{l}', \num{n}))
\lif {} \\
& \qquad (\Obj Q_{q'}(\num{l}, \num{n}') \land \Obj S_{\sigma'}(\num{l}',
\num{n}') \land !A(\num{l}, \num{n}))
\intertext{Caso contrário, seja $l = m = 0$ e considere a seguinte
  !!{sentence} implicada pela segunda fórmula componente:}
& ((\Obj Q_{q_i}(\num{0}, \num{n}) \land \Obj S_{\sigma}(\num{0}, \num{n})) \lif {}\\
& \qquad   (\Obj Q_{q_j}(\num{0}, \num{n}') \land \Obj S_{\sigma'}(\num{0}, \num{n}') \land
!A(\num{0}, \num{n}))) 
\intertext{Qualquer das sentenças implica}
&  \Obj Q_{q'}(\num{l}, \num{n}') \land \Obj S_{\sigma'}(\num{m},
  \num{n}') \land {}\\
&\qquad  \Obj S_{\sigma_0}(\num{0}, \num{n}')\land \dots \land
  \Obj S_{\sigma_k}(\num{k}, \num{n}')  \land {} \\
& \qquad  \lforall[x][(\num{k} < x
  \lif \Obj S_\TMblank(x, \num{n}'))]
\end{align*}
como antes. (Note que no primeiro caso, $\num{l}' \ident \num{l+1}
\ident \num{m}$ e no segundo caso $\num{l} \ident \num{0}$.) Mas
isto é justamente $!C(M, w, n+1)$.

\item O caso \olref{stay} é deixado como exercício.
\end{enumerate}
Mostramos que para qualquer~$n$, $!T(M, w) \Entails !C(M, w, n)$.
\end{proof}

\begin{prob}
Complete o caso~\olref[tur][und][ver]{stay} da prova de
\olref[tur][und][ver]{lem:config}.
\end{prob}

\begin{prob}
Dê !!a{derivation} de $\Obj S_{\sigma_i}(\num{i}, \num{n}')$ a partir de
$\Obj S_{\sigma_i}(\num{i}, \num{n})$ e $!A(m, n)$ (supondo $i \neq
m$, isto é, $i < m$ ou $m < i$).
\end{prob}

\begin{prob}
Dê !!a{derivation} de $\lforall[x][(\num{k}' < x \lif \Obj
  S_\TMblank(x, \num{n}'))]$ a partir de $\lforall[x][(\num{k} < x \lif \Obj
  S_\TMblank(x, \num{n}'))]$, $\lforall[x][x < x']$ e
$\lforall[x][\lforall[y][\lforall[z][ ((x < y \land y < z) \lif x <
      z)]]]$.)
\end{prob}


\begin{lem}
\ollabel{lem:valid-if-halt}
Se $M$ para na entrada~$w$, então $!T(M, w) \lif
!E(M, w)$ é válida.
\end{lem}

\begin{proof}
Por \olref{lem:config}, sabemos que, para qualquer instante~$n$, a
descrição~$!C(M, w, n)$ da configuração de $M$ no instante~$n$ é
implicada por~$!T(M, w)$.  Suponha que $M$ pare após $k$ passos. Nesse
ponto, ela estará examinando o quadrado~$m$, para algum~$m \in \Nat$.
Então $!C(M, w, k)$ descreve uma configuração de parada de~$M$, isto
é, contém como fórmulas componentes tanto $\Obj Q_q(\num{m}, \num{k})$
quanto $\Obj S_\sigma(\num{m}, \num{k})$ com $\delta(q,\sigma)$
indefinida.  Assim, por \olref{lem:halt-config-implies-halt}, $!C(M,
w, k) \Entails !E(M, w)$. Mas como $!T(M, w) \Entails !C(M, w, k)$,
temos $!T(M, w) \Entails !E(M, w)$ e portanto $!T(M, w) \lif !E(M, w)$
é válida.
\end{proof}


\begin{explain} 
Para completar a verificação de nossa afirmação, temos também de
estabelecer a direção reversa: se $!T(M, w) \lif !E(M, w)$ é válida,
então $M$ de fato para quando iniciada na entrada~$w$. 
\end{explain}

\begin{lem}
\ollabel{lem:halt-if-valid}
Se $\Entails !T(M, w) \lif !E(M, w)$, então $M$ para na entrada~$w$.
\end{lem}

\begin{proof}
Considere a $\Lang L_M$-!!{structure}~$\Struct M$ com
domínio~$\Nat$ que interpreta $\Obj 0$ como~$0$, $\prime$~como a função
sucessor e $<$~como a relação menor-que, e os predicados $\Obj Q_q$
e~$\Obj S_\sigma$ como segue:
\begin{align*}
  \Assign{\Obj Q_q}{M} & =
\Setabs{\tuple{m, n}}{\begin{array}{ll}\text{iniciada em~$w$, após~$n$ passos,}\\ \text{$M$ está no estado $q$
  examinando o quadrado~$m$}\end{array}} \\
\Assign{\Obj S_\sigma}{M} & = \Setabs{\tuple{m, n}}{\begin{array}{ll}
\text{iniciada em~$w$, após $n$ passos,}\\ \text{o quadrado~$m$ de $M$ contém
  o símbolo~$\sigma$}\end{array}}
\end{align*}
Em outras palavras, construímos a !!{structure}~$\Struct{M}$ de modo
que ela descreva o que $M$ iniciada na entrada~$w$ de fato faz, passo
a passo. Claramente, $\Sat{M}{!T(M, w)}$. Se $\Entails !T(M, w) \lif !E(M, w)$,
então também $\Sat{M}{!E(M, w)}$, isto é,
\[
\Sat{M}{\lexists[x][\lexists[y][(\bigvee_{\tuple{q, \sigma} \in
      X}(\Obj Q_q(x, y) \land \Obj S_\sigma(x, y)))]]}.
\]
Como $\Domain{M} = \Nat$, deve haver $m$, $n \in \Nat$ tais que
$\Sat{M}{\Obj Q_q(\num{m}, \num{n}) \land \Obj S_\sigma(\num{m},
\num{n})}$ para algum~$q$ e~$\sigma$ tais que $\delta(q, \sigma)$ é
indefinida. Pela definição de~$\Struct M$, isso significa que $M$
iniciada na entrada~$w$ após~$n$ passos está no estado~$q$ e lendo o
símbolo~$\sigma$, e a função de transição é indefinida, isto é,
$M$~parou.
\end{proof}

\end{document}
